``Daksha'' is a proposed Indian mission consisting of two satellites S1 and
S2 orbiting the Earth in the same circular orbit of radius $r = 7000$\,km but
with $180^\circ$ phase difference. These satellites observe the universe in
high energy domain (X-rays and $\gamma$-rays). Each of the satellites of
Daksha uses several flat, rectangular detectors.

To understand how to localize a source in the sky, we shall use a simplified
model of the Daksha mission. Assume that S1 has only two identical detectors
D1 and D2, each of area $A = 0.50\,\mathrm{m}^2$, attached to an opaque mount
M as shown in the figure below. The detectors lie symmetrically around the
$y$ axis in planes perpendicular to the $x$-$y$ plane and make an angle
$\alpha = 120^\circ$ with each other.

\ioaafig{2025_TH_T01-f1.pdf}

\begin{description}
\item[(T01.1)] When observing a distant source located in the $x$-$y$ plane,
  detector D1 records a power $P_1 = 2.70 \times 10^{-10}\,\mathrm{J\,s}^{-1}$
  and detector D2 records a power
  $P_2 = 4.70 \times 10^{-10}\,\mathrm{J\,s}^{-1}$. Estimate the angle $\eta$
  made by the position vector of the source with the positive $y$-axis, with
  counter-clockwise angle from the positive $y$ axis being considered
  positive. \hfill[5]
\end{description}

Consider a single pulse from a distant source (not necessarily in the
$x$-$y$ plane) recorded by both satellites (S1 and S2) of Daksha. The times
of the peaks of the pulses recorded by S1 and S2 are $t_1$ and $t_2$,
respectively.

\begin{description}
\item[(T01.2)] If $t_1 - t_2$ was measured to be $10.0 \pm 0.1$\,ms then
  determine the fraction, $f$, of the celestial sphere where the source
  might lie. \hfill[5]
\end{description}
